\chapter{Dettagli circuitali}
\label{app:circuiti}

I diagrammi sono generati con \code{qml.draw} sui circuiti istanziati dal codice.

\section{Corrispondenza fra teoria e implementazione}

\begin{table}[htbp]
\centering
\caption[Corrispondenza fra formulazione teorica e implementazione]{%
Corrispondenza fra gli elementi teorici del \cref{cap:transformer} e le classi che
li implementano. Salvo diversa indicazione, il file è
\file{experiments/01\_quantum\_gpt/src/model.py}.}
\label{tab:corrispondenza-gpt}
\small
\begin{tabularx}{\textwidth}{@{}l l X@{}}
\toprule
\textbf{Elemento teorico} & \textbf{Classe / metodo} & \textbf{File} \\
\midrule
Embedding di token e posizione & \code{QuantumGPT.\_\_init\_\_} & \file{src/model.py} \\
Attenzione a testa singola & \classe{Head} & idem \\
Maschera causale & \code{Head.tril} (buffer) & idem \\
Attenzione multi-testa & \classe{MultiHeadAttention} & idem \\
Rete feed-forward & \classe{FeedForward} & idem \\
Blocco pre-norm & \classe{Block} & idem \\
Obiettivo autoregressivo & \code{QuantumGPT.forward} & idem \\
Generazione & \code{QuantumGPT.generate} & idem \\
Proiezione $Q,K,V$ quantistica & \classe{QuantumLayerAdapter} & \file{src/quantum\_layers.py} \\
\quad (alias di) & \classe{QuantumLinearWithAdapter} & \file{quantum\_linear.py} \\
\bottomrule
\end{tabularx}
\end{table}

\section{Circuito variazionale degli esperimenti 01--03}

Il circuito è \code{AngleEmbedding} seguito da \code{StronglyEntanglingLayers}
con $\nqlayers = 2$, con misura di $\expval{Z_i}$ su ogni qubit
(\cref{sec:quantum-linear}).

I tre casi differiscono per il numero di fili. Il comando seguente genera
diagramma e specifiche con la versione installata di PennyLane:

\begin{lstlisting}[style=python, numbers=none]
import pennylane as qml
from quantum_framework.layers import QuantumLinear

layer = QuantumLinear(n_qubits=10, n_qlayers=2)
print(qml.draw(layer.qlayer.qnode)(inputs, weights))
print(qml.specs(layer.qlayer.qnode)(inputs, weights))
\end{lstlisting}

\section{Circuito di kernel IQP dell'esperimento 04}

Il circuito \code{kernel\_circuit} realizza
$K(x,y) = \abs{\bra{0}\Uenc^{\dagger}(y)\Uenc(x)\ket{0}}^2$ dell'
\cref{eq:fidelity-kernel} con $\nq = 8$.

Il calcolo applica prima $\Uenc(x)$ e poi $\Uenc^{\dagger}(y)$; la probabilità
di ritrovare $\ket{0}^{\otimes8}$ è la fedeltà. Il controllo a vettore di stato
costruisce esplicitamente i due stati.

\section{Dimensioni dei circuiti}

\begin{center}
\begin{tabular}{@{}l r r r@{}}\toprule
\textbf{Esperimento}&\textbf{Qubit}&\textbf{Pesi VQC}&\textbf{Ampiezze}\\\midrule
01&4&24 per circuito&16\\
02&8&48&256\\
03&10&60&1024\\
04&8&0 addestrabili&256\\\bottomrule
\end{tabular}
\end{center}
Per 01 il modello contiene più istanze del circuito nelle proiezioni $Q,K,V$;
il valore in tabella è riferito a una singola istanza.
